% GENERATED FILE. Edit canonical lessons, then run npm run generate:books.
% canonical-source-hash: fnv1a64:f046dbd35bd75e46
% canonical-lessons: JA-C156-oshiire, JA-C156-hikidashi, JA-C156-shosai, JA-C156-ireru, JA-C156-naraberu

\chapter{Closets, Drawers and a Study}
\label{ch:ja-closets-drawers-and-a-study}
\hlchaptermodality{156}

\begin{chapteropening}
\textbf{By the end of this chapter:} \emph{I can say where things go in a home.}

Complete the last lesson of chapter 156: I can say where things go in a home.
\end{chapteropening}

\section[\texorpdfstring{\ja{おしいれ} (oshiire)}{おしいれ (oshiire)}]{\ja{おしいれ} (oshiire) — a closet (for futons)}

\label{lesson:JA-C156-oshiire}



\begin{quote}
Before the new one: say the Japanese for furniture, then the Japanese for the front entrance.
\end{quote}

\subsection*{You'll want to know: \ja{おしいれ}}
\textbf{\ja{おしいれ}} — \emph{oshiire} — \textquotedblleft{}a closet (for futons)\textquotedblright{}.

\begin{quote}
\textbf{\ja{ふとんを} \ja{おしいれに} \ja{しまいます。}} — \emph{futon o oshiire ni shimaimasu} — \textquotedblleft{}I put the futon away in the closet.\textquotedblright{}
\end{quote}

\begin{cousinweb}
\textbf{\ja{おしいれ}} (\emph{oshiire}) is native Japanese, built from \emph{oshi}, \textquotedblleft{}pushing\textquotedblright{}, the stem of \emph{osu}, and \emph{ire}, \textquotedblleft{}putting in\textquotedblright{}, the stem of \emph{ireru}; both parts are inherited rather than borrowed (JMdict, KANJIDIC2).
\end{cousinweb}

\begin{grammarlens}[title={What you've built}]
One more word for everyday life at A2.
\end{grammarlens}

\subsection*{Your turn}
\begin{itemize}
\raggedright
  \item \emph{Say it:} \emph{oshiire}
  \item \emph{Say it:} \emph{oshiire}, once more
  \item \emph{Say it:} \emph{futon o oshiire ni shimaimasu} — I put the futon away in the closet.
  \item \emph{Recall:} say the Japanese for furniture, then the Japanese for the front entrance, then say \emph{oshiire} again
\end{itemize}

\begin{tcolorbox}[breakable,colback=teal!4,colframe=teal!35!black,title={Before you move on}]
What does \ja{おしいれ} mean? (\textquotedblleft{}A closet (for futons)\textquotedblright{}.) Say it once more.
\end{tcolorbox}
\section[\texorpdfstring{\ja{ひきだし} (hikidashi)}{ひきだし (hikidashi)}]{\ja{ひきだし} (hikidashi) — a drawer}

\label{lesson:JA-C156-hikidashi}



\begin{quote}
Before the new one: say the Japanese for the front entrance, then the Japanese for a closet.
\end{quote}

\subsection*{You'll want to know: \ja{ひきだし}}
\textbf{\ja{ひきだし}} — \emph{hikidashi} — \textquotedblleft{}a drawer\textquotedblright{}.

\begin{quote}
\textbf{\ja{つくえの} \ja{ひきだしは} \ja{ちいさいです。}} — \emph{tsukue no hikidashi wa chiisai desu} — \textquotedblleft{}The desk drawer is small.\textquotedblright{}
\end{quote}

\begin{cousinweb}
\textbf{\ja{ひきだし}} (\emph{hikidashi}) is native Japanese, built from \emph{hiki}, \textquotedblleft{}pulling\textquotedblright{}, the stem of \emph{hiku}, and \emph{dashi}, \textquotedblleft{}taking out\textquotedblright{}, the stem of \emph{dasu}; both parts are inherited rather than borrowed (JMdict, KANJIDIC2).
\end{cousinweb}

\begin{grammarlens}[title={What you've built}]
One more word for everyday life at A2.
\end{grammarlens}

\subsection*{Your turn}
\begin{itemize}
\raggedright
  \item \emph{Say it:} \emph{hikidashi}
  \item \emph{Say it:} \emph{hikidashi}, once more
  \item \emph{Say it:} \emph{tsukue no hikidashi wa chiisai desu} — The desk drawer is small.
  \item \emph{Recall:} say the Japanese for the front entrance, then the Japanese for a closet, then say \emph{hikidashi} again
\end{itemize}

\begin{tcolorbox}[breakable,colback=teal!4,colframe=teal!35!black,title={Before you move on}]
What does \ja{ひきだし} mean? (\textquotedblleft{}A drawer\textquotedblright{}.) Say it once more.
\end{tcolorbox}
\section[\texorpdfstring{\ja{しょさい} (shosai)}{しょさい (shosai)}]{\ja{しょさい} (shosai) — a home office; a study}

\label{lesson:JA-C156-shosai}



\begin{quote}
Before the new one: say the Japanese for a closet, then the Japanese for a drawer.
\end{quote}

\subsection*{You'll want to know: \ja{しょさい}}
\textbf{\ja{しょさい}} — \emph{shosai} — \textquotedblleft{}a home office; a study\textquotedblright{}.

\begin{quote}
\textbf{\ja{ちちは} \ja{しょさいで} \ja{ほんを} \ja{よみます。}} — \emph{chichi wa shosai de hon o yomimasu} — \textquotedblleft{}My father reads books in his study.\textquotedblright{}
\end{quote}

\begin{cousinweb}
\textbf{\ja{しょさい}} (\emph{shosai}) is Sino-Japanese (kango): two Chinese characters read with their on'yomi, readings borrowed from Chinese. They are \emph{sho}, \textquotedblleft{}writing, books\textquotedblright{} (Middle Chinese \emph{syo}) and \emph{sai}, \textquotedblleft{}room, purification\textquotedblright{} (\emph{tsreaj}) (KANJIDIC2; Middle Chinese in Baxter's transcription).
\end{cousinweb}

\begin{grammarlens}[title={What you've built}]
One more word for everyday life at A2.
\end{grammarlens}

\subsection*{Your turn}
\begin{itemize}
\raggedright
  \item \emph{Say it:} \emph{shosai}
  \item \emph{Say it:} \emph{shosai}, once more
  \item \emph{Say it:} \emph{chichi wa shosai de hon o yomimasu} — My father reads books in his study.
  \item \emph{Recall:} say the Japanese for a closet, then the Japanese for a drawer, then say \emph{shosai} again
\end{itemize}

\begin{tcolorbox}[breakable,colback=teal!4,colframe=teal!35!black,title={Before you move on}]
What does \ja{しょさい} mean? (\textquotedblleft{}A home office; a study\textquotedblright{}.) Say it once more.
\end{tcolorbox}
\section[\texorpdfstring{\ja{いれる} (ireru)}{いれる (ireru)}]{\ja{いれる} (ireru) — to put in}

\label{lesson:JA-C156-ireru}



\begin{quote}
Before the new one: say the Japanese for a drawer, then the Japanese for a home office.
\end{quote}

\subsection*{You'll want to know: \ja{いれる}}
\textbf{\ja{いれる}} — \emph{ireru} — \textquotedblleft{}to put in\textquotedblright{}.

\begin{quote}
\textbf{\ja{かばんに} \ja{ほんを} \ja{いれます。}} — \emph{kaban ni hon o iremasu} — \textquotedblleft{}I put a book in my bag.\textquotedblright{}
\end{quote}

Polite \textbf{\ja{いれます}}, \emph{iremasu}; an ichidan verb.

\begin{cousinweb}
\textbf{\ja{いれる}} (\emph{ireru}) is a native Japanese word (wago), inherited rather than borrowed from Chinese, and built on the same root as \emph{iru}, \textquotedblleft{}to go in\textquotedblright{}; its kanji was attached later as a kun'yomi (JMdict, KANJIDIC2).
\end{cousinweb}

\begin{grammarlens}[title={What you've built}]
One more word for everyday life at A2.
\end{grammarlens}

\subsection*{Your turn}
\begin{itemize}
\raggedright
  \item \emph{Say it:} \emph{ireru}
  \item \emph{Say it:} \emph{ireru}, once more
  \item \emph{Say it:} \emph{kaban ni hon o iremasu} — I put a book in my bag.
  \item \emph{Recall:} say the Japanese for a drawer, then the Japanese for a home office, then say \emph{ireru} again
\end{itemize}

\begin{tcolorbox}[breakable,colback=teal!4,colframe=teal!35!black,title={Before you move on}]
What does \ja{いれる} mean? (\textquotedblleft{}To put in\textquotedblright{}.) Say it once more.
\end{tcolorbox}
\section[\texorpdfstring{\ja{ならべる} (naraberu)}{ならべる (naraberu)}]{\ja{ならべる} (naraberu) — to arrange in a row; to line up (things)}

\label{lesson:JA-C156-naraberu}



\begin{quote}
Before the new one: say the Japanese for a home office, then the Japanese for to put in.
\end{quote}

\subsection*{You'll want to know: \ja{ならべる}}
\textbf{\ja{ならべる}} — \emph{naraberu} — \textquotedblleft{}to arrange in a row; to line up (things)\textquotedblright{}.

\begin{quote}
\textbf{\ja{たなに} \ja{ほんを} \ja{ならべます。}} — \emph{tana ni hon o narabemasu} — \textquotedblleft{}I line up the books on the shelf.\textquotedblright{}
\end{quote}

Polite \textbf{\ja{ならべます}}, \emph{narabemasu}; an ichidan verb.

\begin{cousinweb}
\textbf{\ja{ならべる}} (\emph{naraberu}) is a native Japanese word (wago), inherited rather than borrowed from Chinese; its kanji was attached later as a kun'yomi, and its older origin is not settled (JMdict, KANJIDIC2).
\end{cousinweb}

\begin{grammarlens}[title={What you've built}]
One more word for everyday life at A2.
\end{grammarlens}

\subsection*{Your turn}
\begin{itemize}
\raggedright
  \item \emph{Say it:} \emph{naraberu}
  \item \emph{Say it:} \emph{naraberu}, once more
  \item \emph{Say it:} \emph{tana ni hon o narabemasu} — I line up the books on the shelf.
  \item \emph{Recall:} say the Japanese for a home office, then the Japanese for to put in, then say \emph{naraberu} again
\end{itemize}

\begin{tcolorbox}[breakable,colback=teal!4,colframe=teal!35!black,title={Before you move on}]
What does \ja{ならべる} mean? (\textquotedblleft{}To arrange in a row; to line up (things)\textquotedblright{}.) Say it once more.
\end{tcolorbox}
